\section{Introducción}\label{sec:introduccion}

\subsection{El modelo de memoria compartida distribuida}

Un sistema de memoria compartida distribuida reúne la memoria física de varios procesos, que pueden vivir en equipos distintos, y la presenta a las aplicaciones como un único espacio de direcciones. El programa reserva bloques con \texttt{dsm\_alloc}, los recorre con punteros y llama a primitivas de sincronización cuando necesita coordinar el acceso; no hay llamadas de envío ni de recepción en el camino de datos. Lo que comparten los participantes es la dirección virtual y las reglas del protocolo, no un segmento físico común, porque cada proceso conserva su propio objeto de memoria y obtiene el contenido de los demás bajo demanda.

La unidad de coherencia, de transferencia y de protección es la página. Cada proceso proyecta una vista de 4 GiB por omisión en la misma dirección base y la deja inicialmente sin permisos. Cuando la aplicación toca una página que no tiene, el procesador levanta una falta de protección, el manejador de \texttt{SIGSEGV} identifica la página y el tipo de acceso, y un despachador resuelve la solicitud con el nodo que custodia esa página. El contenido se instala por un alias interno y se ajusta el permiso de la vista: lectura para una copia compartida, lectura y escritura para una copia exclusiva. El hilo que sufrió el fallo espera y reintenta la instrucción, así que el programa no distingue un acceso local de uno remoto.

El contenido de una página no vive en un lugar único. Varios procesos pueden tener copias en sus vistas y un nodo custodio, llamado \emph{home}, mantiene la entrada de directorio que indica quién tiene copia, quién la tiene modificada y cuál es la versión vigente. El protocolo MSI garantiza que exista un solo escritor a la vez y que ninguna copia obsoleta sobreviva a una escritura: antes de conceder el estado modificado, el custodio invalida las copias compartidas y espera sus acuses. El costo de esa garantía aparece en el primer acceso a cada página y en las invalidaciones, no en cada lectura posterior.

La coherencia es por página y bajo demanda. Dos variables independientes que comparten página comparten también el costo de la coherencia, y una página que nadie toca no consume red ni estado de directorio más allá de su entrada. El sistema tampoco promete persistencia: la memoria vive mientras viven los procesos que la aportan y el contenido se reconstruye o se pierde según la redundancia configurada.

Desde el punto de vista de la aplicación, el contrato es reducido: inicializar el proceso, reservar bloques, ejecutar con punteros y sincronizar cuando corresponda. La reserva entrega un rango contiguo de la vista y el mismo rango puede compartirse entre procesos cuando uno publica la dirección que obtuvo; los accesos posteriores no llevan tamaño ni identificador, porque el hardware y el manejador ya los describen. Las operaciones colectivas —cerrojos, barreras y kernels distribuidos— se apoyan en la misma vista y en los mismos custodios, así que la lógica del programa no necesita distinguir entre memoria local y remota.

\subsection{Componentes del sistema}

La biblioteca \texttt{libdsm.a} contiene el núcleo completo y se enlaza en cada proceso que participa. Su única cabecera pública, \texttt{dsm.h}, declara el ciclo de vida, la reserva de bloques, las primitivas de sincronización, la ejecución distribuida y las consultas de diagnóstico. El mismo código de coherencia, red y membresía corre en todos los procesos, sin una variante especial para el primero ni para el más antiguo.

Un \emph{nodo} es el proceso que aporta memoria física al espacio global. Se lanza con el binario \texttt{dsm}, lee el archivo de configuración, reserva su bloque de marcos y participa en el descubrimiento, los latidos y el reparto de custodia. Su aporte se declara con \texttt{NODE\_MEM}, que puede variar de un equipo a otro. Un nodo también ejecuta las páginas que le corresponden como custodio y responde las solicitudes de los demás.

Un \emph{cliente} es un proceso que usa la memoria del clúster sin aportar ninguna. Las aplicaciones de prueba se ejecutan en este modo: proyectan la vista, provocan fallos de página y leen o escriben a través del protocolo, pero nunca son custodios ni suman capacidad. El rol se elige con la variable de entorno \texttt{DSM\_NODE\_MODE}, no con una clave del archivo compartido, porque describe al proceso y no al clúster. En la tabla de miembros un cliente aparece en una fila aparte y no altera el conteo de nodos ni la memoria agregada.

Alrededor del núcleo conviven las aplicaciones y las herramientas. \texttt{demo} reúne los escenarios de coherencia, ejecución distribuida y paridad; \texttt{matmul}, \texttt{juego\_vida} y \texttt{filtro\_video} ejercitan la API sobre cargas de trabajo reconocibles; \texttt{dsmctl} consulta un nodo en ejecución sin unirse al clúster y \texttt{video\_pack} prepara material de entrada. El reporte HTML, generado por un nodo o por la herramienta, resume la membresía, el directorio, las páginas materializadas y las franjas de paridad.

La biblioteca también expone contadores y consultas de diagnóstico: fallos de lectura y escritura, mensajes y bytes, páginas en cada estado, transiciones de latidos, deltas de paridad y franjas reconstruidas. \texttt{dsmctl} los presenta desde fuera del clúster y el reporte HTML los organiza por nodo, con el custodio, el copyset y la franja de cada página materializada. La observación ocurre por una conexión efímera de solo lectura, así que diagnosticar un clúster en marcha no altera su membresía ni su estado de coherencia.

\subsection{Panorama de la arquitectura}

El sistema es simétrico y no tiene coordinador de coherencia. Cada nodo toma sus decisiones con la tabla de membresía, la época vigente y el anillo de custodia, tres estructuras que todos calculan con las mismas reglas. La tabla se replica mediante los saludos y los anuncios; el anillo se deriva del identificador y de la memoria aportada por cada miembro, por lo que dos nodos que ven la misma membresía llegan al mismo custodio para una página sin consultarse.

Sobre esa base se distinguen tres planos. El plano de datos transporta páginas, concesiones, invalidaciones y acuses según el estado MSI. El plano de control mantiene la membresía: descubrimiento por anuncios, unión por \texttt{JOIN}, ranuras, épocas y latidos con sospecha. El plano de servicios ofrece reserva global de bloques, cerrojos y barreras distribuidos, ejecución de kernels por particiones y reportes de diagnóstico. Los tres comparten una cabecera serializada de 56 bytes sobre TCP y un canal UDP para el descubrimiento.

Un nodo que arranca no conoce de antemano a los demás. Anuncia su nombre, dirección, puerto, rol y memoria por el grupo multicast \texttt{239.255.42.99:27699}, con broadcast como respaldo, y escucha los anuncios ajenos. El archivo \texttt{dsm.conf} describe el clúster y no una lista de equipos: las claves estructurales deben coincidir y se resumen en un hash que viaja en el saludo, mientras la memoria aportada y los parámetros locales quedan fuera de esa comparación. La sección~\ref{sec:arquitectura} desarrolla los componentes y el ciclo de vida; la sección~\ref{sec:configuracion} detalla el archivo, el descubrimiento y la membresía.

\subsection{Organización del documento}

Después de esta introducción, la arquitectura general describe la biblioteca, los roles de nodo y cliente, los grupos de hilos, el despliegue físico, el ciclo de vida y la organización del código fuente. La configuración, el descubrimiento y la membresía explican el archivo compartido, el hash de compatibilidad, los anuncios, la asignación de ranuras, los latidos, las épocas y la heterogeneidad de memoria.

Los capítulos siguientes recorren el camino de un acceso. La memoria y la vista virtual presentan el objeto compartido, los marcos y los permisos por página; la coherencia MSI describe los estados, el directorio y las transacciones; la concurrencia y el manejo de señales explican cómo el manejador de fallos convive con los despachadores, los trabajadores y la red. La redundancia y la paridad cubren las franjas, los deltas y la reconstrucción; la ejecución distribuida explica el registro de kernels y el reparto por particiones.

El protocolo de mensajería documenta la cabecera, los tipos y las reglas de validación; la API pública reúne las firmas y los códigos de estado; las aplicaciones de prueba y las herramientas de diagnóstico muestran el sistema en uso. La validación reúne las pruebas de extremo a extremo, las mediciones del clúster local y la ejecución sobre equipos físicos. Las conclusiones cierran el recorrido y los apéndices conservan el material de referencia que no conviene intercalar en la lectura principal.
